% answers_short.tex —— 2020 年·全国Ⅲ卷·理科 · 简短答案速查
% 由 tools/make_short_answers.py 自 answers.tex 的 \gksolhead/\gkanswer 提取，
% 并与 2020/2020_数学真题_全国3理_答案要点.md 逐题比对后生成，勿手改（改源 MD 后重跑）。

\begin{gknotice}
\textbf{使用说明}：本节只给最终答案便于快速核对，完整推导见解析版。答案为编写者独立推导，\textbf{非官方标准答案}。
\end{gknotice}

\gkpart{选择题}
\begin{center}\begin{tabular}{ccccccccccccc}
\toprule
题号 & 1 & 2 & 3 & 4 & 5 & 6 & 7 & 8 & 9 & 10 & 11 & 12 \\
\midrule
答案 & C & D & B & C & B & D & A & C & D & D & A & A \\
\bottomrule
\end{tabular}\end{center}

\gkpart{填空题}
\begin{itemize}[leftmargin=2.6em]
  \item[\textbf{13.}] $7$。可行域是顶点为 $(0,0)$，$(1,-1)$，$(1,2)$ 的三角形，$z=3x+2y$ 在 $(1,2)$ 处取最大值 $7$。
  \item[\textbf{14.}] $240$。通项 $\mathrm C_6^k2^kx^{12-3k}$ 中令 $12-3k=0$ 得 $k=4$，常数项为 $\mathrm C_6^4\cdot2^4=240$。
  \item[\textbf{15.}] $\dfrac{\sqrt2}3\pi$。圆锥高 $2\sqrt2$，轴截面内切圆半径 $\rho=\dfrac{1\times2\sqrt2}{1+3}=\dfrac{\sqrt2}2$，体积 $\dfrac43\pi\rho^3=\dfrac{\sqrt2}3\pi$。
  \item[\textbf{16.}] ②③。$f$ 为奇函数，图象关于原点对称；又 $f(\pi-x)=f(x)$，图象关于直线 $x=\dfrac\pi2$ 对称；而 $\sin x<0$ 时 $f(x)$ 无下界，①④ 均假。
\end{itemize}

\gkpart{解答题}
\begin{itemize}[leftmargin=2.6em]
  \item[\textbf{17.}] (1) $a_2=5$，$a_3=7$，$a_n=2n+1$。由递推式算得前两项，再用数学归纳法验证 $3(2k+1)-4k=2(k+1)+1$；(2) $S_n=(2n-1)2^{n+1}+2$。错位相减求和，或以 $n=1,2,3$ 验证
  \item[\textbf{18.}] (1) $0.43$，$0.27$，$0.21$，$0.09$。四个等级的天数依次为 $43$，$27$，$21$，$9$，各除以 $100$；(2) $350$。$\dfrac{20\times100+35\times300+45\times500}{100}=350$；(3) $K^2\approx5.820>3.841$，有 $95\%$ 的把握认为有关（列联表见解析版）
  \item[\textbf{19.}] (1) 证明。取 $A$ 为原点建系，得 $\overrightarrow{AC_1}=\overrightarrow{AE}+\overrightarrow{AF}$，故 $C_1$ 在平面 $AEF$ 内；(2) $\dfrac{\sqrt{42}}7$。两平面法向量可取 $(1,1,-1)$ 与 $(-1,-4,-2)$，其夹角余弦为 $-\dfrac1{\sqrt7}$，故二面角的正弦值为 $\dfrac{\sqrt{42}}7$
  \item[\textbf{20.}] (1) $\dfrac{x^2}{25}+\dfrac{16y^2}{25}=1$。由 $e^2=\dfrac{25-m^2}{25}=\dfrac{15}{16}$ 得 $m^2=\dfrac{25}{16}$；(2) $\dfrac52$。设 $Q(6,q)$，由垂直且等长得 $P$ 为 $(5-q,1)$ 或 $(5+q,-1)$，代回椭圆得 $q=\pm2$ 或 $q=\pm8$，四种情形的面积同为 $\dfrac52$
  \item[\textbf{21.}] (1) $b=-\dfrac34$。切线与 $y$ 轴垂直给出 $f'\left(\dfrac12\right)=\dfrac34+b=0$；(2) 证明。设 $|x_0|\leqslant1$ 为零点，则 $f(x)=(x-x_0)\left(x^2+x_0x+x_0^2-\dfrac34\right)$，其余零点满足 $\left|x\right|\leqslant\dfrac{|x_0|+\sqrt{3-3x_0^2}}2\leqslant1$
\end{itemize}

\gkpart{选做题}
\begin{itemize}[leftmargin=2.6em]
  \item[\textbf{22.}] (1) $4\sqrt{10}$。令 $y=0$ 得 $t=2$、$A(-4,0)$，令 $x=0$ 得 $t=-2$、$B(0,12)$，故 $|AB|=\sqrt{16+144}=4\sqrt{10}$；(2) $3\rho\cos\theta-\rho\sin\theta+12=0$（即 $\rho(\sin\theta-3\cos\theta)=12$）。直线 $AB$ 的直角坐标方程为 $3x-y+12=0$，代入 $x=\rho\cos\theta$，$y=\rho\sin\theta$ 即得
  \item[\textbf{23.}] (1) 证明。由 $(a+b+c)^2=0$ 得 $ab+bc+ca=-\dfrac{a^2+b^2+c^2}2$，而 $abc=1$ 保证 $a^2+b^2+c^2>0$，故 $ab+bc+ca<0$；(2) 证明。设 $a$ 为最大值，则 $a>0$ 且 $b$，$c$ 为 $t^2+at+\dfrac1a=0$ 的两实根，判别式 $a^2-\dfrac4a\geqslant0$，故 $a^3\geqslant4$，即 $a\geqslant\sqrt[3]{4}$
\end{itemize}
